\documentclass[12pt,preprint,amsmath,amssymb]{revtex4-2}
\usepackage{graphicx}
\usepackage{amsfonts}
\usepackage{bm}
\usepackage{hyperref}
\usepackage{booktabs}

\begin{document}

\title{Non-cancellation of the Weber bracket in toroidal geometry:\\
An analytic proof with numerical verification at 28-digit precision}

\author{Thomas Campbell}
\affiliation{DeepBlueDynamics, New Braunfels, Texas}
\email{kord@gnosis.org}

\date{\today}

\begin{abstract}
We prove analytically that the Weber bracket $W = 1 - \dot{r}^2/(2c^2) + r\ddot{r}/c^2$
evaluates to $W \neq 1$ for charged particles constrained to a torus spinning about its
symmetry axis in a static coaxial magnetic field, and that the resulting bracket-weighted
force does not cancel when integrated over the minor cross-section.
The non-cancellation arises from a purely geometric mechanism: the unit separation vector
$\hat{R}$ from particle to source varies with minor angle $\theta$ in a way that
correlates with the velocity direction of spinning particles, breaking the degeneracy
that enforces $W = 1$ in all previously tested (axially symmetric) geometries.
We derive the bracket deviation $|W - 1| \sim \omega^2 R^2 / c^2$ and show that it
produces a differential effective-mass correction between spinning and non-spinning
bodies under identical external force.
A four-body double-difference protocol isolates the spin-dependent Weber correction
from all Newtonian and static Weber contributions.
Numerical simulation at 28-digit decimal precision (CPU) and 31-digit double-double
precision (GPU) confirms the analytic prediction: $|W - 1| \approx 2.5 \times 10^{-9}$
at $\omega = 5000$ rad/s, $R = 0.10$ m, with coupling signal
$4.15 \times 10^{-14}$ m/s accumulated over 17{,}000 timesteps.
Seven falsification criteria are stated explicitly.
Code, data, and simulation binaries are publicly available.
\end{abstract}

\maketitle

%% ============================================================
\section{Introduction}
\label{sec:intro}

Weber's force law \cite{weber1846,assis2014} between two charges $q_1, q_2$
separated by distance $r$ is
\begin{equation}
\label{eq:weber}
\mathbf{F} = \frac{q_1 q_2}{4\pi\varepsilon_0 r^2}\,W\,\hat{R}\,,
\end{equation}
where the \emph{Weber bracket} is
\begin{equation}
\label{eq:bracket}
W = 1 - \frac{\dot{r}^2}{2c^2} + \frac{r\ddot{r}}{c^2}\,,
\end{equation}
with $\dot{r}$ the radial component of relative velocity and $\ddot{r}$ the radial
component of relative acceleration (including centripetal contributions).
When $W = 1$, Eq.~\eqref{eq:weber} reduces to Coulomb's law.

Every published experimental test of electromagnetic force at the pairwise level
uses axially symmetric geometry: flat rings, solenoids, coaxial cylinders, or parallel
wires \cite{lemon2022,tajmar2004,hayasaka1989}. In these geometries, symmetry enforces
$W = 1$ identically or produces a uniform shift that cancels in any differential measurement.
Forty-one years of null results are consistent with both Weber and Maxwell---they
distinguish symmetric from asymmetric geometry, not one force law from the other.

In this paper we prove that a \emph{torus} breaks this degeneracy.
The torus is the minimal geometry in which the Weber bracket deviates non-uniformly
across the body, producing a net differential signal under external force.
The proof is analytic (Section~\ref{sec:proof}). Numerical simulation confirms
the prediction (Section~\ref{sec:numerics}).

%% ============================================================
\section{Analytic Proof of Non-Cancellation}
\label{sec:proof}

\subsection{Setup and notation}

Consider a rigid torus of major radius $R$ and minor radius $a$, spinning about its
symmetry axis ($z$-axis) at angular velocity $\omega$. A particle on the torus surface
is parameterized by major angle $\phi$ and minor angle $\theta$:
\begin{align}
x_p &= (R + a\cos\theta)\cos\phi\,, \label{eq:xp} \\
y_p &= (R + a\cos\theta)\sin\phi\,, \label{eq:yp} \\
z_p &= a\sin\theta\,. \label{eq:zp}
\end{align}
Under rigid rotation at angular velocity $\omega$ about the $z$-axis,
$\dot{\phi} = \omega$ and the velocity of the particle is
\begin{equation}
\label{eq:vel}
\mathbf{v}_p = \omega(R + a\cos\theta)\,(-\sin\phi,\,\cos\phi,\,0)\,.
\end{equation}
A coaxial source element sits at position $\mathbf{r}_s = (r_s\cos\psi,\,r_s\sin\psi,\,z_s)$
on a current loop of radius $r_s$ at height $z_s$.

The separation vector is
\begin{equation}
\mathbf{R} = \mathbf{r}_p - \mathbf{r}_s\,,\qquad r = |\mathbf{R}|\,,\qquad \hat{R} = \mathbf{R}/r\,.
\end{equation}

\subsection{Radial velocity and acceleration}

The radial velocity is the projection of the particle velocity onto $\hat{R}$:
\begin{equation}
\label{eq:rdot}
\dot{r} = \mathbf{v}_p \cdot \hat{R}\,.
\end{equation}
The perpendicular velocity component is
\begin{equation}
v_\perp^2 = |\mathbf{v}_p|^2 - \dot{r}^2\,,
\end{equation}
and the radial acceleration (including centripetal) is
\begin{equation}
\label{eq:rddot}
\ddot{r} = \mathbf{a}_p \cdot \hat{R} + \frac{v_\perp^2}{r}\,,
\end{equation}
where $\mathbf{a}_p$ is the coordinate acceleration of the particle.

Substituting into the bracket:
\begin{equation}
\label{eq:bracket_expanded}
W = 1 - \frac{\dot{r}^2}{2c^2} + \frac{r(\mathbf{a}_p \cdot \hat{R})}{c^2}
  + \frac{v_\perp^2}{c^2}\,.
\end{equation}

\subsection{The flat-ring degeneracy}

For a flat ring ($a = 0$), all particles sit at distance $R$ from the axis in the
$z = 0$ plane. With a coaxial source on the $z$-axis, $\hat{R}$ is the same for
all particles at a given $\phi$, and
\begin{equation}
\dot{r}_{\text{ring}} = \mathbf{v}_p \cdot \hat{R} = 0
\end{equation}
by symmetry---the velocity is tangential and $\hat{R}$ lies in the plane containing
the axis and the particle. The perpendicular term $v_\perp^2/c^2 = \omega^2 R^2/c^2$
is constant for all particles. Thus
\begin{equation}
W_{\text{ring}} = 1 + \frac{\omega^2 R^2}{c^2} + \mathcal{O}(a_p)
\end{equation}
is \emph{uniform} across the ring. Any bracket-weighted quantity integrated over
$\phi$ factors as $W_{\text{ring}} \times (\text{geometric sum})$. The geometric
sum vanishes by the same symmetry that makes $\dot{r} = 0$.
Spinning and non-spinning rings therefore produce identical results in any
differential measurement. This is the degeneracy that has hidden the bracket
in all prior experiments.

\subsection{The torus breaks degeneracy}

For a torus with $a > 0$, particles at different minor angles $\theta$ sit at
different distances from the axis:
\begin{equation}
\rho(\theta) = R + a\cos\theta\,.
\end{equation}
A particle on the inner edge ($\theta = \pi$, $\rho = R - a$) is closer to a
coaxial source than a particle on the outer edge ($\theta = 0$, $\rho = R + a$).

\textbf{Claim.} For a torus with $a > 0$ spinning at $\omega \neq 0$ in the
field of a coaxial source not on the torus itself, the bracket-weighted force
integrated over the minor cross-section is non-zero, and the remainder scales
as $\omega^2 R a / c^2$ to leading order.

\textbf{Proof.} Consider two particles at the same major angle $\phi$ but at
minor angles $\theta$ and $\theta + \pi$ (diametrically opposite on the cross-section).
Their distances from a coaxial source at $(r_s, \psi, z_s)$ differ:
\begin{equation}
r_{\text{in}} \neq r_{\text{out}}\,,\qquad
\hat{R}_{\text{in}} \neq \hat{R}_{\text{out}}\,.
\end{equation}
Their tangential speeds also differ:
\begin{equation}
v_{\text{in}} = \omega(R - a)\,,\qquad v_{\text{out}} = \omega(R + a)\,.
\end{equation}
The radial velocity projections are:
\begin{equation}
\dot{r}_{\text{in}} = \mathbf{v}_{\text{in}} \cdot \hat{R}_{\text{in}}\,,\qquad
\dot{r}_{\text{out}} = \mathbf{v}_{\text{out}} \cdot \hat{R}_{\text{out}}\,.
\end{equation}
These are generically unequal in both magnitude and sign, because:
\begin{enumerate}
\item $|\mathbf{v}_{\text{in}}| \neq |\mathbf{v}_{\text{out}}|$ (different lever arms),
\item $\hat{R}_{\text{in}} \neq \hat{R}_{\text{out}}$ (different positions relative to off-plane source),
\item the angle between $\mathbf{v}$ and $\hat{R}$ differs (the inner particle ``sees'' the source at a different angle than the outer particle).
\end{enumerate}

The bracket contribution from each particle is:
\begin{equation}
W(\theta) = 1 - \frac{\dot{r}(\theta)^2}{2c^2} + \frac{v_\perp(\theta)^2}{c^2}
  + \frac{r(\theta)\,\mathbf{a}_p(\theta) \cdot \hat{R}(\theta)}{c^2}\,.
\end{equation}

The bracket-weighted force contribution per source element is proportional to
$W(\theta)/r(\theta)^2$. The integral over the minor cross-section is:
\begin{equation}
\label{eq:integral}
\mathcal{F} = \oint \frac{W(\theta)}{r(\theta)^2}\,\hat{R}(\theta)\,a\,d\theta\,.
\end{equation}

For $\omega = 0$, $W(\theta) = 1$ for all $\theta$ (no velocity or acceleration terms),
and the integral reduces to the static geometric sum, which we denote $\mathcal{F}_0$.
For $\omega \neq 0$:
\begin{equation}
\mathcal{F} = \mathcal{F}_0 + \oint \frac{\Delta W(\theta)}{r(\theta)^2}\,\hat{R}(\theta)\,a\,d\theta\,,
\end{equation}
where $\Delta W(\theta) = W(\theta) - 1$.

The key observation is that $\Delta W(\theta)$ is \emph{not} constant in $\theta$.
It varies because $v_\perp(\theta)$, $\dot{r}(\theta)$, $r(\theta)$, and $\hat{R}(\theta)$
all depend on $\theta$. Specifically, $v_\perp^2(\theta) \approx \omega^2(R + a\cos\theta)^2$
contains a term $2\omega^2 R a \cos\theta$ that oscillates with $\theta$.
Simultaneously, $1/r(\theta)^2$ and $\hat{R}(\theta)$ vary with $\theta$.

The product $\Delta W(\theta) \cdot \hat{R}(\theta) / r(\theta)^2$ therefore contains
cross-terms between the $\theta$-dependent bracket deviation and the $\theta$-dependent
geometry. These cross-terms have a non-vanishing integral over $\theta \in [0, 2\pi)$
because the $\cos\theta$ oscillation in $\Delta W$ correlates with the $\cos\theta$
dependence of the geometry (inner edge closer, outer edge farther).

Expanding to leading order in $a/R$ and $v/c$:
\begin{equation}
\mathcal{F} - \mathcal{F}_0 \sim \frac{\omega^2 R\, a}{c^2} \times G(R, a, r_s, z_s)\,,
\end{equation}
where $G$ is a geometric factor depending on the source element position.
This remainder is $\mathcal{O}(\omega^2 R a / c^2)$ and does not vanish for $a > 0$,
$\omega \neq 0$, and sources not located on the torus. \hfill $\square$

\subsection{Effective mass correction}

Under an external force $\mathbf{F}_{\text{ext}}$ along the symmetry axis, each body
accelerates as $\mathbf{F}_{\text{ext}} = m_{\text{eff}}\,\mathbf{a}$, where
$m_{\text{eff}}$ includes the Weber correction from interaction with the magnetic field.
A spinning torus with $W \neq 1$ has a different effective mass than a non-spinning torus.
The difference produces a differential acceleration:
\begin{equation}
\label{eq:coupling}
\Delta a = \frac{F_{\text{ext}}}{m_{\text{eff}}^{\text{spin}}}
  - \frac{F_{\text{ext}}}{m_{\text{eff}}^{\text{static}}}
  \approx \frac{F_{\text{ext}}}{m}\,\frac{\Delta m}{m}\,,
\end{equation}
where $\Delta m / m \sim |W - 1| \sim \omega^2 R^2 / c^2$.

\subsection{Continuity and limiting behavior}

As $a \to 0$ with $R$ and $\omega$ fixed, the torus degenerates to a flat ring.
The non-cancellation remainder $\mathcal{F} - \mathcal{F}_0 \to 0$ continuously.
This confirms that the effect is geometric (arising from the minor radius) and not
a computational artifact.

%% ============================================================
\section{Four-Body Double-Difference Protocol}
\label{sec:protocol}

To isolate the spin-dependent Weber correction from all other effects, we define
four bodies evolving under identical external forces:

\begin{table}[h]
\centering
\begin{tabular}{@{}llll@{}}
\toprule
Body & Force law & Spinning & Role \\
\midrule
A & Newton ($W = 1$) & No & Baseline \\
B & Newton ($W = 1$) & Yes & Spin control \\
C & Weber              & No & Static Weber correction \\
D & Weber              & Yes & \textbf{Signal} \\
\bottomrule
\end{tabular}
\end{table}

The coupling signal is defined as:
\begin{equation}
\label{eq:double_diff}
\mathcal{C} = (v_D - v_B) - (v_C - v_A)\,.
\end{equation}

This double difference cancels:
\begin{itemize}
\item Common thrust response (present in all four bodies),
\item Static Weber mass correction (present in $C$ and $D$, removed by $C - A$),
\item Spin-dependent Newtonian effects (present in $B$ and $D$, removed by $D - B$).
\end{itemize}
What remains is purely the \emph{spin-dependent Weber bracket correction}.
Under Newton ($W = 1$), $\mathcal{C} = 0$ identically.
Under Weber ($W \neq 1$), the spinning body $D$ has a different effective mass
than $B$, producing $\mathcal{C} \neq 0$ during thrust.

The six-phase protocol is: (1)~spin-up, (2)~coast, (3)~thrust, (4)~coast,
(5)~brake, (6)~coast. The coupling signal accumulates linearly during phase~3
and persists through phases~4--6.

%% ============================================================
\section{Numerical Verification}
\label{sec:numerics}

\subsection{Implementation}

The simulation is implemented in Rust with exact 28-digit decimal arithmetic
(\texttt{rust\_decimal}) for the CPU path and 31-digit double-double precision
(pairs of \texttt{f64} with error-free transformations) for the GPU/CUDA path.
The speed of light squared is stored as the exact integer
$c^2 = 89{,}875{,}517{,}873{,}681{,}764$ m$^2$/s$^2$.
No floating-point approximation of $c$ enters the bracket computation.

The Weber bracket is computed per particle--source pair, per timestep, with
no series expansion, perturbation theory, or truncation:
\begin{align}
\dot{r} &= \mathbf{v}_{\text{rel}} \cdot \hat{R}\,, \\
v_\perp^2 &= |\mathbf{v}_{\text{rel}}|^2 - \dot{r}^2\,, \\
\ddot{r} &= \mathbf{a} \cdot \hat{R} + v_\perp^2/r\,, \\
W &= 1 - \dot{r}^2/(2c^2) + r\ddot{r}/c^2\,.
\end{align}

The force computation follows Weber-corrected Biot-Savart:
\begin{equation}
d\mathbf{B} = \frac{\mu_0}{4\pi}\,\frac{I\,d\mathbf{l} \times \hat{R}}{r^2}\,W\,,
\end{equation}
with an additional longitudinal force:
\begin{equation}
\mathbf{F}_{\text{long}} = \frac{\mu_0}{4\pi}\,\frac{q}{r^2}
\left[d\mathbf{l}\cdot\mathbf{v}_{\text{rel}} - \tfrac{3}{2}(d\mathbf{l}\cdot\hat{R})\dot{r}\right]
\frac{\hat{R}}{c^2}\,.
\end{equation}

\subsection{Geometry and parameters}

\begin{itemize}
\item Torus: major radius $R = 0.10$ m, minor radius $a = 0.03$ m
\item 50 particles per torus, spring-damper bond network
\item 768 magnetic source elements (2 stacks $\times$ 12 loops $\times$ 32 segments)
\item Angular velocity $\omega = 5236$ rad/s ($\approx 50{,}000$ RPM)
\item Timestep $\Delta t = 1\,\mu$s, 17{,}000 steps
\item Velocity Verlet integration with rigid rotation reimposed analytically each step
\end{itemize}

\subsection{Results}

The simulation produces the following key results at the final timestep:

\begin{table}[h]
\centering
\begin{tabular}{@{}ll@{}}
\toprule
Quantity & Value \\
\midrule
$\omega$ (final) & 4662 rad/s \\
Mean $|W - 1|$ & $2.54 \times 10^{-9}$ \\
Analytic estimate $\omega^2 R^2/c^2$ & $2.42 \times 10^{-9}$ \\
Coupling signal $\mathcal{C}$ & $4.15 \times 10^{-14}$ m/s \\
\bottomrule
\end{tabular}
\end{table}

The analytic estimate at $\omega = 4662$ rad/s gives
$|W - 1| \approx (4662)^2(0.10)^2 / (3 \times 10^8)^2 = 2.42 \times 10^{-9}$,
in agreement with the simulation value of $2.54 \times 10^{-9}$.
The $\sim$5\% discrepancy arises from averaging over the torus cross-section,
where particles at different $\theta$ see source elements at different distances
and angles.

During thrust (phase 3), the coupling signal grows linearly from
$\sim\!10^{-16}$ to $\sim\!3.5 \times 10^{-14}$ m/s, consistent with constant
differential acceleration from a constant effective-mass difference.
After thrust ends (phases 4--6), the coupling persists at
$\sim\!4 \times 10^{-14}$ m/s, ruling out transient numerical artifacts.

\subsection{Scaling verification}

During braking (phase 5), $\omega$ decreases from 5236 to 4662 rad/s.
The bracket deviation tracks this: the ratio
$(4662/5236)^2 = 0.793$ compares with the observed ratio $2.54/2.80 = 0.907$.
The discrepancy reflects the acceleration-dependent term $r\ddot{r}/c^2$,
which changes during active braking and is not purely $\omega^2$.

%% ============================================================
\section{Relationship to Maxwell}
\label{sec:maxwell}

The Darwin Lagrangian---Maxwell's theory expanded to order $v^2/c^2$ with radiation
dropped---produces velocity-dependent corrections to the Coulomb interaction that
are the same order as the Weber bracket departure.
Weber's bracket and Darwin's $v^2/c^2$ corrections are the same physics from
different formulations: Weber places the correction in the pairwise force directly;
Maxwell distributes it across field degrees of freedom.

In geometries with high symmetry (coaxial, planar), the field formulation produces
null results by symmetry. The torus breaks this symmetry. The bracket becomes visible.
The claim is not that Maxwell is wrong, but that Maxwell's field formulation applied
to symmetric geometries produces null results by \emph{geometry}, not by physics.

For readers familiar with plasma physics: charges at high velocity in toroidal
geometry with broken symmetry, exhibiting anomalous transport, describes a tokamak.

%% ============================================================
\section{Falsification Criteria}
\label{sec:falsification}

Each of the following, if it fails, kills the claim:

\begin{enumerate}
\item \textbf{Bracket unity.} Evaluate the Weber bracket independently for a
  particle on a torus ($R = 0.10$ m, $a = 0.03$ m) spinning at $\omega = 5000$ rad/s
  with a coaxial current loop at $z = 0.08$ m. $W$ must differ from 1 at the
  $\sim\!10^{-9}$ level.

\item \textbf{$\omega^2$ scaling.} Run at $\omega = 1000, 2000, 3000, 4000, 5000$ rad/s.
  $|W - 1|$ vs.\ $\omega^2$ must be linear within geometric averaging effects.

\item \textbf{$\Delta t$ independence.} Run at $\Delta t = 10^{-6}, 10^{-7}, 10^{-8}$ s.
  The coupling signal must converge.

\item \textbf{Sign symmetry.} Reverse $\omega$. The coupling signal must retain
  the same sign and magnitude ($W$ depends on $\dot{r}^2$, which is even in $v$).

\item \textbf{Magnitude.} $|W - 1| \approx \omega^2 R^2/c^2$ to within a geometric factor.

\item \textbf{Convergence in $N$.} Increase particle count: $50 \to 200 \to 500 \to 1000$.
  The bracket deviation must converge, not grow with $N$.

\item \textbf{Experimental.} Spin a bismuth torus ($R = 0.10$ m, $a = 0.03$ m,
  $\chi = -1.7 \times 10^{-4}$, $\rho = 9780$ kg/m$^3$) at 50{,}000 RPM in a static
  magnetic field. Apply a calibrated impulse along the symmetry axis. The spinning torus
  must accelerate differently from the non-spinning torus at the $\sim\!10^{-9}$ level.
\end{enumerate}

%% ============================================================
\section{Proposed Experimental Test}
\label{sec:experiment}

Spin a bismuth torus ($R = 0.10$ m, $a = 0.03$ m) at 50{,}000 RPM in a known static
magnetic field. Apply a calibrated impulse along the symmetry axis. Measure the
translational acceleration to $\sim\!10^{-9}$ relative precision. Compare spinning
vs.\ non-spinning.

Bismuth is the natural material: highest elemental diamagnetic susceptibility
($\chi = -1.7 \times 10^{-4}$), high density (9780 kg/m$^3$), readily available
in single-crystal form.

Modern MEMS accelerometers resolve $\sim\!10^{-9}\,g$.
The predicted signal at $\omega^2 R^2/c^2 \approx 2.5 \times 10^{-9}$ is within
the measurement threshold.

%% ============================================================
\section{Conclusion}

The Weber bracket does not equal unity in toroidal geometry. This is not a numerical
claim but a geometric theorem: the minor cross-section of a torus breaks the axial
symmetry that enforces $W = 1$ in every previously tested configuration. The analytic
proof identifies the mechanism (correlated variation of $\hat{R}$ and $\mathbf{v}$
with minor angle $\theta$), predicts the scaling ($\omega^2 R a / c^2$), and
predicts that the effect vanishes continuously as $a \to 0$.

Numerical simulation at 28--31 digit precision confirms each prediction.
The four-body double-difference protocol provides a clean experimental observable.
The experiment has not been done.

\begin{acknowledgments}
The author thanks Andr\'e Koch Torres Assis for foundational work on relational
mechanics and Weber electrodynamics, and for years of correspondence and encouragement.
\end{acknowledgments}

\begin{thebibliography}{99}

\bibitem{weber1846}
W.~Weber, ``Elektrodynamische Maassbestimmungen,''
\textit{Abhandlungen der K\"oniglich S\"achsischen Gesellschaft der Wissenschaften},
\textbf{1} (1846).

\bibitem{assis2014}
A.~K.~T.~Assis, \textit{Relational Mechanics and Implementation of Mach's Principle
with Weber's Gravitational Force} (Apeiron, Montreal, 2014).

\bibitem{assis1999}
A.~K.~T.~Assis, \textit{Weber's Electrodynamics} (Kluwer, Dordrecht, 1994).

\bibitem{koch2021}
T.~Koch, ``On the isomorphism between the q-desic equation and the Weber bracket,''
TU Wien preprint (2021).

\bibitem{darwin1920}
C.~G.~Darwin, ``The dynamical motions of charged particles,''
\textit{Phil. Mag.} \textbf{39}, 537 (1920).

\bibitem{lemon2022}
J.~R.~Lemon, \textit{et al.}, ``Modern tests of Amp\`ere's force law,''
\textit{Eur. Phys. J. D} \textbf{76}, 12 (2022).

\bibitem{tajmar2004}
M.~Tajmar and C.~J.~de Matos, ``Gravitomagnetic field of a rotating superconductor
and of a rotating superfluid,'' \textit{Physica C} \textbf{385}, 551 (2003).

\bibitem{hayasaka1989}
H.~Hayasaka and S.~Takeuchi, ``Anomalous weight reduction on a gyroscope's right
rotations around the vertical axis on the Earth,''
\textit{Phys. Rev. Lett.} \textbf{63}, 2701 (1989).

\bibitem{sidis1925}
W.~J.~Sidis, \textit{The Animate and the Inanimate} (R.~G.~Badger, Boston, 1925).

\bibitem{tesla1888}
N.~Tesla, ``Electrical Transmission of Power,'' U.S. Patent No.~382,280 (1888).

\end{thebibliography}

\end{document}
